\subsection{Macaulay modules over a local ring}

PROPOSITION 3. - Let A be a Noetherian local ring, M a nonzero finitely generated A-module, and d the dimension of M. The following conditions are equivalent:

\begin{enumerate}
\def\labelenumi{(\roman{enumi})}
\item
  the A-module M is Macaulay
\item
  we have \(\operatorname{prof}(\mathbf{M}) = d\) ;
\item
  we have \(\mathrm{Ext}_{\Lambda}^{i}(\kappa_{\mathrm{A}},\mathrm{M}) = 0\) for every integer \(i < d\) ;
\item
  we have \(\mathrm{Ext}_{\mathrm{A}}^{i}(\mathrm{N},\mathrm{M}) = 0\) for every A-module N of finite length and every integer \(i < d\) ;
\item
  we have \(\operatorname{Ext}_A^i (\mathrm{N},\mathrm{M}) = 0\) for every finitely generated A-module N and every integer \(i < d - \dim (\mathrm{M}\otimes_{\mathrm{A}}\mathrm{N})\) ;
\item
  there exists an M-regular sequence of elements of \(\mathfrak{m}_{\mathrm{A}}\) of length \(d\) .
\end{enumerate}

Condition (i) is equivalent to the equality prof(M) = d, that is, to condition (ii), or again to the inequality prof(M) ≥ d, that is, to (iii) and (vi) (§ 1, no. 4, th. 2). The implications (v) ⇒ (iv) and (iv) ⇒ (iii) are evident.

Finally, suppose M is Macaulay and let N be a finitely generated A-module. Set \(F = \text{Supp}(N)\); by II, § 4, no. 4, prop. 18, we have \(\text{Supp}(M) \cap F = \text{Supp}(M \otimes_{A} N)\) , so that

\[
\operatorname{prof} _ {F} (M) = \operatorname{codim} (\operatorname{Supp} (M) \cap F, \operatorname{Supp} (M)) = \dim M - \dim (M \otimes_ {A} N)
\]

(No. 1, cor. of Prop. 1 and No. 2, cor. of Prop. 2). The implication (i)⇒(v) then follows from Remark 1 of § 1, No. 5.

In the rest of this section, we shall say that a finitely generated module M over a Noetherian local ring A is pure if, for every prime ideal p associated with M, we have \(\dim(\mathrm{A}/\mathfrak{p}) = \dim(\mathrm{M})\). This also means that M has no embedded associated prime ideals and that the irreducible components of the support of M all have the same dimension. Every Macaulay module over a Noetherian local ring is pure (no. 2, prop. 2 and its corollary).

Lemma 1.-- Let A be a Noetherian local ring, M a finitely generated and pure A-module, and x an element of \(\mathfrak{m}_{\mathrm{A}}\). The following conditions are equivalent:

\begin{enumerate}
\def\labelenumi{(\roman{enumi})}
\item
  we have \(\dim(M/xM)=\dim(M)-1\) ;
\item
  the homothety \(x_{M}\) is injective.
\end{enumerate}

We may assume M is nonzero. Assertion (i) is equivalent to the fact that x belongs to none of the minimal elements p of Supp(M) such that dim(A/p) = dim(M) (VIII, § 3, no. 2, prop. 3), and assertion (ii) is equivalent to the fact that x belongs to none of the prime ideals associated with M (IV, § 1, no. 1, cor. 2 of prop. 2). Since M is pure, its associated prime ideals are the minimal elements of Supp(M), hence (i) and (ii) are equivalent.

Let A be a Noetherian local ring and M a nonzero finitely generated A-module. Recall (VIII, § 3, no. 2) that a sequence \((x_{1},\ldots,x_{r})\) of elements of \(m_{A}\) is said to be secant for M if one has \(\dim(M/(x_{1}M+\ldots+x_{r}M))=\dim(M)-r\) .

PROPOSITION 4. Let A be a Noetherian local ring, M a nonzero finitely generated A-module, and \((x_{1},\ldots ,x_{r})\) be a sequence of elements of \(\mathfrak{m}_{A}\) secant for M. The following conditions are equivalent:

\begin{enumerate}
\def\labelenumi{(\roman{enumi})}
\item
  the A-module M is Macaulay
\item
  the sequence \((x_{1},\ldots,x_{r})\) is M-regular and the A-module \(\mathrm{M}/(x_{1}\mathrm{M}+\ldots+x_{r}\mathrm{M})\) is Macaulay.
\end{enumerate}

Suppose that the sequence \((x_{1},\ldots,x_{r})\) is M-regular. We then have

\[
\dim (\mathrm{M}) = r + \dim (\mathrm{M} / (x _ {1} \mathrm{M} + \dots + x _ {r} \mathrm{M}))
\]

\[
\operatorname{prof} (\mathrm{M}) = r + \operatorname{prof} \left(\mathrm{M} / (x _ {1} \mathrm{M} + \dots + x _ {r} \mathrm{M})\right)
\]

(§ 1, no. 4, prop. 7), whence the implication (ii)⇒(i).

Suppose the A-module M is Macaulay and let us prove (ii) by induction on \(r\). The assertion is evident if \(r = 0\). If \(r \geqslant 1\) , the A-module \(N = M/(x_1M + \ldots + x_{r-1}M)\) is Macaulay by the induction hypothesis and one has \(\dim(N/x_rN) = \dim(N) - 1\) since the sequence \((x_1, \ldots, x_r)\) is secant; the homothety \((x_r)_N\) is therefore injective (lemma 1), and \(N/x_rN\) is Macaulay (\(\mathrm{no.}\) 1, example 3), whence (ii).

THEOREM 1.--- Let \(A\) be a Noetherian local ring, M a nonzero finitely generated A-module, d the dimension of M, and let \(\mathbf{x} = (x_1, \ldots, x_d)\) be a sequence of elements of \(\mathfrak{m}_{\mathrm{A}}\) secant for M, and J the ideal it generates. The following conditions are equivalent:

\begin{enumerate}
\def\labelenumi{(\roman{enumi})}
\item
  the A-module M is Macaulay
\item
  the sequence x is M-regular;
\item
  the sequence x is completely secant for M (A, X, p.~157, def. 2);
\item
  the multiplicity (VIII, § 7, no. 1, def. 1) \(e_{\mathrm{J}}(\mathrm{M})\) of M relative to the ideal J is equal to the length of the A-module M/JM ;
\item
  for every integer i such that \(1 \leqslant i \leqslant d\) , the A-module \(\mathrm{M}/(x_{1}\mathrm{M}+\ldots+x_{i-1}\mathrm{M})\) is pure.
\end{enumerate}

The equivalence of (ii) and (iii) follows from A, X, p.~160, cor. 1 of th. 1. Since the A-module M/JM has finite length (VIII, § 3, no. 2, th. 1), the equivalence of (iii) and (iv) follows from VIII, § 4, no. 3, prop. 4 and no. 4, th. 3. It remains to prove the equivalence of (i), (ii), and (v).

(i)⇒(v): if M is Macaulay, each of the modules \(\mathrm{M}/(x_{1}\mathrm{M}+\ldots+x_{i-1}\mathrm{M})\) is Macaulay (prop. 4), hence pure.

(v)⇒(ii): this follows from Lemma 1 applied to each of the modules M/(x₁M+\ldots+\(x_{i-1}\){}M).

(ii)⇒(i): this follows from Prop. 4, since M/JM is of finite length, hence Macaulay.

